\documentclass[11pt]{article}
\usepackage[margin=0.81in]{geometry}
\usepackage{amsmath,amssymb,amsthm}
\usepackage{titling}
\setlength{\droptitle}{-24pt}
\posttitle{\par\end{center}\vskip -10pt}
\postdate{\par\end{center}\vskip -16pt}
\usepackage[hidelinks]{hyperref}
\newtheorem{theorem}{Theorem}
\newtheorem{lemma}[theorem]{Lemma}
\newtheorem{corollary}[theorem]{Corollary}
\newcommand{\Z}{\mathbb Z}
\newcommand{\F}{\mathbb F}
\title{Two generators for every finite metabelian quotient\protect\\of an explicit prime knot family}
\author{}
\date{18 November 2026}
\begin{document}\maketitle
\begin{abstract}
Every finite metabelian quotient of the knot group in an explicit
prime four-bridge twist family is generated by at most two elements.
We prove this by showing that the Alexander module needs at most one
generator over every residue field of the Laurent ring. Two fixed
Fox rows reduce the calculation to one specialization in characteristic
$83$. A linear combination of two minors handles that specialization
uniformly in the twist parameter.
\end{abstract}

Let $G_n$ denote the knot group of the closure of
\[
\sigma_1\sigma_2\sigma_3^{-2}\sigma_2\sigma_3^{-2}\sigma_2
\sigma_1^2\sigma_2^3\sigma_3^{-1}\sigma_2^{3+6n},\qquad n\ge0.
\]
The companion proof \cite[Lemma~1]{fox} gives the four-meridian
presentation with closure relators $R_i=F_ix_i^{-1}$ for $0\le i\le3$,
where $F_i$ is the free word in bottom slot $i$. The first three
relators imply the fourth by ordered-product preservation. We keep
all four rows for the calculation. For each braid letter the positive
rule is $(a,b)\mapsto(b,b^{-1}ab)$ and the negative rule is
$(a,b)\mapsto(aba^{-1},a)$.

\begin{theorem}\label{thm:metabelian}
For every $n\ge0$, every finite metabelian quotient of $G_n$ can be
generated by at most two elements. Every nonabelian quotient in this class
has ordinary rank two.
\end{theorem}

Write $\Lambda=\Z[t,t^{-1}]$, and let $H_{ji}$ be the image of
$\partial R_j/\partial x_i$ under the map $x_i\mapsto t$.
The matrix $H$ has rows indexed by the four closure relators.
The theorem follows from a unit ideal of $2$-by-$2$ minors of $H$.

\begin{lemma}\label{lem:minors}
For every $n\ge0$, the $2$-by-$2$ minors of $H$ generate $\Lambda$.
\end{lemma}
\begin{proof}
Write $h_i$ for row $i$. The terminal twists change only slots $1,2$,
so $R_0,R_3$ are each fixed as free words. Compute the two Laurent
minors on rows $(0,3)$ and columns $(0,3)$ and $(1,3)$, and multiply
both by $t^5$. The resulting polynomials are
\begin{align*}
 f(t)&=t^9-3t^8+4t^7-8t^5+18t^4-21t^3+15t^2-6t+1,\\
 g(t)&=-2t^8+9t^7-22t^6+36t^5-45t^4+40t^3-25t^2+10t-2.
\end{align*}
The Sylvester determinant is $\operatorname{Res}(f,g)=-83$. Modulo
$83$ the monic greatest common divisor of $f,g$ is $t+24$.
These exact calculations use free-word Fox differentiation, a
$17$-by-$17$ integer determinant and Euclidean division over $\F_{83}$.
The companion code \texttt{verify\_scalar\_cocycle\_capacity.py}
replays them.

Suppose a field specialization $t\mapsto u\ne0$ kills every
$2$-by-$2$ minor. Then $f(u)=g(u)=0$. Evaluation at $u$ gives a
nonzero vector in the kernel of the Sylvester matrix, so the field
has characteristic $83$. The greatest common divisor then forces
$u=59$ in its prime field.

At this specialization the first row and the second row at $n=0,1$
are
\[
 h_0=(10,12,60,1),\qquad h_1(0)=(51,66,7,42),\qquad
 h_1(1)=(1,26,55,1).
\]
Let $m_{01}(n),m_{02}(n)$ be the two minors on rows $(0,1)$ and
columns $(0,1),(0,2)$. At $n=0$ their values are $(48,81)$ and at
$n=1$ they are $(82,75)$.

Before subtracting the identity rows to form the closure derivatives,
the pair of middle bottom derivative rows transforms at each terminal
crossing by the $2$-by-$2$ matrix
\[
 B(u)=\begin{pmatrix}0&1\\ u^{-1}&1-u^{-1}\end{pmatrix}.
\]
Its eigenvalues are $1$ and $-u^{-1}$. At $u=59$ they are distinct,
and the second eigenvalue of $B(u)^6$ is $81\ne1$. Thus $h_1(n)$
is an affine linear function of $81^n$ over $\F_{83}$, and the
minors involving the fixed row $h_0$ have the same property.
The function
\[
                   77m_{01}(n)-34m_{02}(n)
\]
takes the value $29$ at both $n=0$ and $n=1$. Since $1$ and $81$ are
distinct, it equals $29$ for every $n$. Therefore the two minors cannot
both vanish at this specialization.

If the ideal of minors were proper, reduction modulo a maximal ideal
containing it would give a field specialization at which they all
vanish. The preceding argument excludes every such specialization.
\end{proof}

\begin{proof}[Proof of Theorem~\ref{thm:metabelian}]
Let $M=G_n'/G_n''$ be the Alexander module, with $t$ acting by
conjugation by $x_0$. To compute its presentation, take the infinite
cyclic cover of the presentation complex. Its vertices are indexed
by the integers. Each generator edge advances by one vertex, so the
boundary on the free edge module $\Lambda^4$ is
$(t-1,t-1,t-1,t-1)$. The kernel consists of vectors with coordinate
sum zero and is free of rank three. The lifted relator boundaries
are the rows of $H$, since every letter contributes the oriented edge
at its current vertex. Therefore $M$ is the quotient of the
rank-three kernel by those rows. This also follows by abelianizing
the fundamental group of the cyclic cover, which is $G_n'$.
For any residue field $k$ of $\Lambda$, the rank-three kernel is
identified with the coordinate-sum-zero subspace of $k^4$ using its
basis $e_i-e_0$ for $i=1,2,3$. Lemma~\ref{lem:minors} gives at least
two independent relation rows in that space. Hence
$\dim_k(M\otimes_\Lambda k)\le1$.

Now let $Q$ be a finite metabelian quotient of $G_n$ and let $A=Q'$.
The image of $x_0$ generates $Q/A$, and its conjugation action makes
$A$ a finite quotient of the $\Lambda$-module $M$. If $A=0$ then
$Q$ is cyclic. Otherwise the image $S$ of $\Lambda$ in
$\operatorname{End}(A)$ is a finite commutative ring with identity.
For each maximal ideal $\mathfrak m$ of $S$, the field
$S/\mathfrak m$ is a residue field of $\Lambda$. The preceding bound
implies $\dim_{S/\mathfrak m} A/\mathfrak m A\le1$.

For each $\mathfrak m$ choose an element of $A$ generating
$A/\mathfrak m A$. Use the Chinese remainder theorem to take a single
linear combination $a\in A$ with these residues at all maximal ideals.
Then $C=A/Sa$ satisfies $C/\mathfrak m C=0$ for every $\mathfrak m$.
If the finite module $C$ were nonzero, it would have a simple
quotient $S/\mathfrak m$ for some maximal ideal, contradicting this
vanishing. Thus $A=Sa$. The conjugates of $a$ by powers of the image
of $x_0$ generate $A$, so that meridian image and $a$ generate $Q$.
A nonabelian group requires at least two generators.
\end{proof}

\begin{corollary}\label{cor:scalar}
For any field $k$ and scalar action $x_i\mapsto u\in k^\times$,
the space of crossed homomorphisms $G_n\to(k,+)$ has dimension at most two.
\end{corollary}
\begin{proof}
Values on the four meridians identify crossed homomorphisms
with the kernel of the specialized $H$. Indeed the cocycle product
rule gives exactly the Fox relation equations. The kernel has
dimension at most two by Lemma~\ref{lem:minors}.
\end{proof}

The ordinary rank of $G_n$ remains open.

\paragraph{AI assistance and review.}
ChatGPT (OpenAI) assisted with mathematical arguments, literature checks,
exposition, proof review and verification code.

\begin{thebibliography}{9}
\bibitem{fox} \emph{Uniform Fox-ideal model barriers for a prime
four-bridge family}. Companion proof, 18 November 2026.
\end{thebibliography}
\end{document}
